\section{Discussion}
\label{sec:discussion}

\subsection{What the model has learned}

The Fourier analysis in Section \ref{sec:analysis} establishes that the
multi-task model uses two separate bases inside one residual stream: an
additive-group Fourier basis that handles addition and subtraction, and
a multiplicative-group Fourier basis that handles multiplication. This
is the more parsimonious of the two outcomes we hypothesised in Section
\ref{sec:intro}: rather than discovering a single ``algebraic'' feature
set that abstracts over group structure, the model has effectively
time-multiplexed the two known per-task algorithms.

This is consistent with the architectural prior. Different cyclic groups
require different Fourier bases (the bases live on index sets of
different sizes), so a single basis that worked for both groups would
have to be either over-complete or coarser than each task individually
needs. Under aggressive weight decay, the implicit bias of the
optimiser is toward Frobenius-economical representations, and economy
in this setting means specialisation.

\subsection{Curriculum vs joint training}

The curriculum experiment (training on addition first, then introducing
subtraction and multiplication) is not faster than joint training from
scratch on a per-step basis. The intuition is that the addition-only
warm-start fills the residual stream with additive-group features at
the cost of capacity that is later needed for the multiplicative basis;
the network spends additional steps undoing this commitment when
multiplication is introduced.

\subsection{Limitations}

Three limitations are worth noting:
\begin{itemize}
\item The analysis is restricted to one prime ($p = 113$). The
      robustness sweep (Section \ref{sec:multitask}) checks two
      additional primes for the addition baseline, but we did not run
      the full multi-task analysis at multiple primes.
\item All experiments use a one-layer transformer with four attention
      heads, matching \citet{nanda2023progress}. A two-layer model would
      have a richer mechanistic story but a less unambiguous
      interpretability picture.
\item The cosine-similarity feature-overlap metric in Figure
      \ref{fig:fig8_feature_overlap} is a coarse summary --- it tells
      us that the same frequencies are weighted in similar proportions,
      not that the embeddings are pointwise equal in those subspaces.
      A finer-grained analysis would compute the principal angles
      between the relevant subspaces.
\end{itemize}

\subsection{Relation to prior work}

Our results extend the single-task mechanistic story of
\citet{nanda2023progress} to the multi-task setting and confirm that the
algorithmic content of grokking is robust to the introduction of
additional related tasks. They are consistent with the effective theory
of \citet{liu2022understanding}: weight decay pushes the network
toward the simplest representation that solves the training set, and
in the multi-task case ``simplest'' turns out to be two separate
bases rather than a single shared one.

The slingshot mechanism of \citet{thilak2022slingshot} predicts that
adaptive optimisers can produce sudden generalisation transitions for
reasons unrelated to weight decay. We do not see evidence of this
mechanism dominating in the present experiments --- the dynamics are
slow and steady at the published $\beta_2 = 0.98$ --- but a careful
disentanglement at smaller $\beta_2$ would be a useful follow-up.
